\section{Results}\label{sec:results}

The two headline empirical results of the project are an open-loop
benchmark on the canonical $84$-day horizon (one chronic time constant
of $42$ days, doubled) and the closed-loop fourteen-day rolling
controller. Both compare the SMC${}^2$ schedule to the canonical
constant baseline $\Phi \equiv 1.0$. Headline numbers are summarised in
Table~\ref{tab:results}.

\begin{table}[h]
\centering
\small
\begin{tabular}{@{}lccccc@{}}
\toprule
Benchmark & Horizon & $\bar A$ (SMC${}^2$) & $\bar A$ (baseline) & gain & F-violation \\
\midrule
Open-loop, Stage D     & $84$ d & $0.645$  & $0.503$  & $+28\%$ & $1.28\%$ \\
Closed-loop MPC, Stage E5 & $14$ d & $0.0964$ & $0.0823$ & $+17\%$ & $0.00\%$ \\
\bottomrule
\end{tabular}
\caption{Headline empirical results. $\bar A$ denotes the mean
autonomic amplitude $\frac{1}{T}\int_0^T A(t)\,\dd t$ averaged over Monte
Carlo trajectories. Sources:
\texttt{version\_1/outputs/fsa\_high\_res/RESULT.md} (Stage D),
\texttt{version\_2/outputs/fsa\_high\_res/RESULT.md} (Stage E5).}
\label{tab:results}
\end{table}

\subsection{Stage D: open-loop schedule on the $84$-day horizon}

In the open-loop benchmark, the dynamics parameters $\thetadyn$ are
treated as known at their truth values, the initial state is fixed at
$(B_0,F_0,A_0) = (0.05,0.30,0.10)$ (a de-trained athlete with high
residual fatigue), and the SMC${}^2$ control of Section~\ref{sec:control}
is run \emph{once} over the full $84$-day horizon. The output is a single
posterior-mean schedule $\Phi(\,\cdot\,;\bar\theta)$.

The optimised schedule has mean $\Phi \approx 1.96$ over the horizon and
a front-loaded periodisation profile (heavy training for the first
$\sim 60$ days, taper thereafter). Mean amplitude
$\bar A = 0.645$ exceeds the constant baseline $\bar A = 0.503$ by
$28\%$, exceeds the sedentary reference $\bar A = 0.209$ by a factor of
$3.08$, and the soft fatigue barrier is violated only $1.28\%$ of the
time --- well within the $5\%$ acceptance gate of
Section~\ref{sec:mpc-feasibility}. All four acceptance gates (mean-A
matches baseline within $3\%$, mean-A $\ge 1.40\times$ sedentary, mean-$\Phi
\in [0.5,2.5]$, F-violation $\le 5\%$) pass.

\subsection{Stage E5: closed-loop MPC on the $14$-day macrocycle}

In the closed-loop benchmark the dynamics parameters are no longer known
to the controller. Algorithm~\ref{alg:fsa-mpc} is run on a
$14$-day record of four-channel observations, with a one-day filter
window, twelve-hour stride, and $27$ replans. The plant is the
\texttt{StepwisePlant} of Section~\ref{sec:mpc-code}.

Mean amplitude under the closed-loop policy is $\bar A = 0.0964$, against
$\bar A = 0.0823$ for the constant baseline $\Phi \equiv 1.0$ on the same
$14$-day horizon. (The shorter horizon explains the smaller absolute
numbers compared to Stage D: the integral $\int_0^T A(t)\,\dd t$ has not
yet converged to the asymptotic mean.) The $+17\%$ gain is achieved with
a remarkable schedule shape: the controller commits a
\emph{rest-leaning} $\Phi$ in the range $0.25$--$0.31$ across all $27$
strides, against the canonical baseline of $1.0$. Physiologically, this
is the controller exploiting the fact that the simulated athlete starts
with high residual fatigue ($F_0 = 0.30$): the optimal short-horizon
strategy is to clear fatigue rather than to add training stimulus, which
unlocks the Stuart--Landau growth-rate $\mu(B,F)$ and lets $A$ rise from
$0.10$ to $\sim 0.20$ over the macrocycle.

The fatigue ceiling is never violated ($0.00\%$). Wall-clock cost on a
single GPU is $41$ minutes for the full closed loop. The closed-loop
policy is robust despite imperfect filter coverage --- on this record
only $3$ of $26$ estimable parameters meet a strict $\ge 5/6$ coverage
gate at the window level --- consistent with the flat-cost-surface
argument of Section~\ref{sec:mpc-robustness}.

\medskip
A trajectory plot of $A(t)$ under both policies, and a plot of the
SMC${}^2$ posterior-mean schedule, are included as Figure
placeholders here; they can be regenerated from the artefacts in
\texttt{version\_2/outputs/fsa\_high\_res/} by re-running
\texttt{tools/bench\_smc\_full\_mpc\_fsa.py}.

\begin{center}
\fbox{\parbox{0.85\textwidth}{\centering
\textit{Figure placeholder: $A(t)$ trajectory under closed-loop MPC
(blue) and constant baseline $\Phi \equiv 1$ (grey), with the
discovered rest-leaning $\Phi$ schedule overlaid.}}}
\end{center}
